% Hoja Cornell con cuadrícula (A4) — Eric Lucero González — https://ericlucero.dev
% Licencia CC BY-NC-SA 4.0. Compilar con pdfLaTeX (dos pasadas: latexmk lo hace solo).
\documentclass[a4paper]{article}
\usepackage[T1]{fontenc}
\usepackage[spanish]{babel}
\usepackage{lmodern}
\usepackage{tikz}
\pagestyle{empty}

% ── Configuración (mm) ─────────────────────────────────────
% Papel A4: 210 x 297 mm. Con estos valores el área de notas mide
% 130 x 210 mm: un número exacto de cuadros de 5 mm.
\def\anchoPapel{210}
\def\altoPapel{297}
\def\margen{10}        % margen exterior del marco
\def\cabecera{17}      % franja superior (título y fecha)
\def\columna{60}       % columna izquierda (ideas clave, preguntas)
\def\resumen{50}       % franja inferior (resumen)
\def\paso{5}           % lado de cada cuadro de la cuadrícula

% Etiquetas de cada zona: \etiquetasfalse para una hoja limpia.
\newif\ifetiquetas
\etiquetastrue

\begin{document}
\begin{tikzpicture}[
    remember picture, overlay, shift={(current page.south west)},
    x=1mm, y=1mm,
    grid lines/.style={gray!45, very thin},
    structure lines/.style={black!75, thick},
    label/.style={gray!70, font=\sffamily\scriptsize, anchor=north west, inner sep=1.5mm}
  ]
  % Coordenadas de las cuatro zonas
  \pgfmathsetmacro{\xIzq}{\margen}
  \pgfmathsetmacro{\xDer}{\anchoPapel - \margen}
  \pgfmathsetmacro{\yAbajo}{\margen}
  \pgfmathsetmacro{\yArriba}{\altoPapel - \margen}
  \pgfmathsetmacro{\xNotas}{\margen + \columna}
  \pgfmathsetmacro{\yNotas}{\margen + \resumen}
  \pgfmathsetmacro{\yCabecera}{\altoPapel - \margen - \cabecera}

  % Cuadrícula del área de notas
  \draw[step=\paso, grid lines, shift={(\xNotas,\yNotas)}]
    (0,0) grid (\xDer-\xNotas, \yCabecera-\yNotas);

  % Marco exterior
  \draw[structure lines] (\xIzq,\yAbajo) rectangle (\xDer,\yArriba);
  % Cabecera (título / fecha)
  \draw[structure lines] (\xIzq,\yCabecera) -- (\xDer,\yCabecera);
  % Columna izquierda (ideas clave / preguntas)
  \draw[structure lines] (\xNotas,\yNotas) -- (\xNotas,\yCabecera);
  % Resumen
  \draw[structure lines] (\xIzq,\yNotas) -- (\xDer,\yNotas);

  \ifetiquetas
    \node[label] at (\xIzq,\yArriba)   {Título};
    \node[label] at (\xDer-40,\yArriba) {Fecha};
    \node[label] at (\xIzq,\yCabecera) {Ideas clave};
    \node[label] at (\xIzq,\yNotas)    {Resumen};
  \fi
\end{tikzpicture}
\end{document}
